% Katalog F: Abbildungen (Palette Slot 6).
\ZSFNewColumn
\StartChapter[kat:bilder]{F · Bilder}

\begin{runintext}
  Drei Makros für drei Fälle: die eigenständige Abbildung, Bild neben Text
  und das Bild ohne Box. Alle nehmen \ZSFkeyword*{width} und
  \ZSFkeyword*{height} selbst entgegen; jeder andere Schlüssel geht an die
  umgebende \ZSFkeyword{figbox}.
\end{runintext}

\SubsectionBar[kat:bild-fig]{Eigenständige Abbildung}
\ZSFindex{ZSFfig}
\ZSFindex{ZSFfigside}

\ZSFfig[width=0.5]{example-image-a}{F-01 · ZSFfig}%
  {Mustertext zum Formvergleich. Ohne width und height greifen die zentralen
   Vorbelegungen.}

\ZSFfig{example-image-a,example-image-b,example-image-c}%
  {F-02 · ZSFfig mit Pfadliste}%
  {Dasselbe Makro, mehrere Pfade: eine Reihe gleich breiter Bilder, die
   Breite folgt der Anzahl.}

\ZSFfigside[width=0.3]{example-image-a}{F-03 · ZSFfigside}{%
  Mustertext zum Formvergleich. Der erste Wert bestimmt den Bildanteil, der
  Rest gehört dem Text.}

\SubsectionBar[kat:bild-image]{Bild ohne eigene Box}
\ZSFindex{ZSFimage}

\begin{figbox}[F-04 · ZSFimage in der figbox][align=center]
  \ZSFimage[width=0.55]{example-image-b}
  \ZSFfigcaption{F-05 · ZSFfigcaption — nur wenn eine figbox direkt befüllt
  wird; die Bild-Makros setzen ihre Beschriftung selbst.}
\end{figbox}
\Zweck{Dasselbe Makro wie in der Tabellenzelle (D-09), hier mit dem
Höhenbudget einer Abbildung. Die Höhe stellt der Container.}

\SubsectionBar[kat:bild-tikz]{Selbstgezeichnetes Diagramm}
\ZSFindex{figbox}

\begin{figbox}[F-06 · figbox mit tikzpicture]
  \centering
  \begin{tikzpicture}
    \begin{axis}[
      width=\linewidth,
      height=\ZSFfigMaxHeight,
      axis lines=middle,
      xmin=-2.2, xmax=2.2,
      ymin=-0.3, ymax=4.5,
      samples=40,
      ticks=none
    ]
      \addplot[color=\chaptercolor, very thick, domain=-2:2]{x^2};
      \node[\chaptercolor, font=\ZSFfontDiagramLabel, anchor=center]
        at (rel axis cs:0.72,0.78) {Mustertext};
    \end{axis}
  \end{tikzpicture}
  \ZSFfigcaption{Der Fall, für den die figbox gedacht ist: kein Bild aus
  einer Datei, sondern eine Zeichnung im Dokument.}
\end{figbox}
\Zweck{Braucht das Opt-in-Modul 12\_plots für die axis-Umgebung; tikzpicture
allein genügt ohne es.}
